\documentclass[../DoAn.tex]{subfiles}
\begin{document}

\begin{center}
    \Large{\textbf{ABSTRACT}}\\
\end{center}
\vspace{1cm}
% Interim update 2026-10-28. TODO: update once final E1/E2 results with a local LLM and observed-data results are available.
\textbf{Background.} Microscopic simulation is the standard tool for analysing signalised intersections, but the default behaviour models of SUMO assume lane-disciplined, rule-compliant traffic. Mixed car and motorcycle traffic in Hanoi breaks these assumptions: vehicles run red lights in swarms, start before the green, weave laterally, stop in the right-turn lane, ride on the sidewalk in queues, and turn left illegally by waiting inside the junction box. Tuning parameters by hand to reproduce these behaviours is laborious and hard to trace.

\textbf{Objective.} This thesis reproduces road-user behaviour at the Le Van Luong -- Hoang Dao Thuy intersection (Thanh Xuan, Hanoi) during the morning peak, 07:30--09:00, as closely as possible. The scenario contains 11 behaviours in three groups: signal non-compliance, gap seeking in congestion, and interaction inside the junction. The illegal left turn with waiting in the box is one of the 11.

\textbf{Method.} A locally hosted large language model (LLM) generates and iteratively refines a driver-behaviour rule specification. The specification holds static vehicle-type parameters and conditional rules executed through TraCI. The LLM produces rules outside the simulation loop; SUMO and a TraCI controller execute them. Each specification is validated against a schema and a source-documented parameter space before it is compiled and simulated.

\textbf{Results so far.} The whole workflow is implemented: a reproducible network, the validator, the compiler, a controller with five conditional rules, the M1--M10 metrics with an aggregate score, an evaluation harness with a shared simulation budget, the LLM loop, and an Optuna baseline. Ten of the 11 behaviours are in the model; wrong-way riding is deferred after a feasibility experiment. The local LLM is \texttt{qwen3.5:9b} (Q4\_K\_M quantisation) served by Ollama. The Optuna baseline has used its full budget of 36 runs; none of its 12 specifications meets the eligibility criteria, because the base model is overloaded at the proxy demand. The two LLM campaigns (parameter tuning on a fixed structure, and generation of structure and parameters) have no final results yet, so the LLM and Optuna cannot be compared.

\textbf{Limitations.} No video-based observations exist yet, so the workflow runs on proxy days derived from published studies of other intersections. Results on these data show that the workflow operates; they do not show how closely the simulation matches real traffic. \textbf{[TODO: final E1/E2 results with a local LLM]}

\end{document}
